\documentclass[twocolumn]{aastex631}

\usepackage{amsmath}
\usepackage{multirow}
\usepackage{natbib}
\usepackage{graphicx} 
\usepackage{aas_macros}

\begin{document}

Many scalar-tensor theories within the Horndeski class, despite their rich phenomenology, face significant challenges with quantum instabilities, such as ghost degrees of freedom, and often struggle with observational viability. This study proposes that disformal transformations act as a fundamental selection principle, identifying robust Horndeski sub-classes that possess emergent non-linear spacetime symmetries, which inherently protect them from these quantum pathologies. Through analytical derivations of disformal invariance conditions, identification of emergent symmetries, a one-loop perturbative analysis, and verification against cosmological and gravitational wave observations, we demonstrate this principle. Focusing on the Galileon theory, a class deeply connected to disformal structures, we show that disformal transformations are crucial for generating its emergent non-linear spacetime symmetry. This Galilean symmetry, rigorously confirmed via a one-loop perturbative analysis utilizing Ward identities, ensures a non-renormalization of the scalar kinetic term, thus protecting the theory from ghost instabilities. Furthermore, this quantum-stable class naturally predicts the speed of gravitational waves to be unity, consistent with GW170817, provides viable late-time cosmic acceleration (validated by numerical background evolution), and offers distinct, testable signatures for large-scale structure growth, including a significant deviation in the gravitational slip parameter (derived from linear perturbation analysis). This work establishes a profound link between disformal structures, emergent symmetries, quantum stability, and observational viability, offering a robust framework for physics beyond General Relativity.

\end{document}
                